\documentclass[11pt]{article}

\usepackage[margin=1in]{geometry}
\usepackage{array}
\usepackage{longtable}
\usepackage{hyperref}

\setlength{\emergencystretch}{3em}

\title{OmegaSim Results Summary and External Feedback Brief}
\author{Prepared for Ben Goertzel}
\date{2026-06-28}

\begin{document}
\maketitle

\begin{abstract}
OmegaSim has now produced a useful, mostly conservative body of evidence.
The successful results are not yet evidence for structured strange attractors,
but they are scientifically valuable: the simulator harness is reproducible;
task pressure, service capacity, exogenous demand, and two-hive transfer
coupling all have measurable and interpretable effects; bounded predictors
can improve forecast skill under matched demand streams; and source-ledger
and null-control infrastructure can now fail closed when apparent structure is
explained by accounting variables. The main null result is equally important:
the completed A2--A7 experiments do not yet justify claims of lobe grammar,
semantic recurrence, strange-attractor-like dynamics, or causal collective
structure. The active next direction is A7.2: delayed artifact-mediated
endogenous prediction in a single hive, followed regardless of outcome by a
separately preregistered three-hive ring experiment family.
\end{abstract}

\section{Purpose and Current Status}

This document summarizes OmegaSim results to date for external feedback on
next scientific directions. It distinguishes three categories:

\begin{enumerate}
\item successful simulator and mechanism findings;
\item null or fail-closed results that constrain interpretation;
\item next experiments aimed at eliciting richer nonlinear collective
dynamics.
\end{enumerate}

The current repository state is clean and synchronized with the remote at the
A7.2 config-contract gate. The latest source-of-truth status says:

\begin{quote}
Run A7.2-style delayed artifact-mediated endogenous prediction first; after
that, whether positive or negative, move to a separately preregistered
three-hive ring.
\end{quote}

The A7.2 preregistration and config/schema contract are already written. Full
test verification has passed, including the focused A7.2 contract tests and
the full harness test suite. No result-bearing A7.2 simulations have yet been
interpreted.

\section{Simulator Abstraction}

OmegaSim is a deterministic abstract simulator for OmegaHive-like collective
cognition. It does not call live LLMs, theorem provers, browsers, Slack,
Atomspace, task boards, or external services. Instead, those are represented
numerically through agents, queues, task classes, actions, graph connections,
artifact fields, transfer events, budgets, and metrics.

The baseline single-hive model has 15 agents divided among coordinator,
researcher, architect, implementer, and reviewer roles. Agents operate over a
task queue and choose actions such as idling, messaging, task creation, and
work. Later experiments add attention policies, exogenous arrivals, service
capacity variation, multi-hive transfer, predictive control, logistic appraisal
state, semantic/artifact fields, and source-ledger accounting.

The most important methodological lesson is that raw lobe labels are not
sufficient evidence. Early labels such as backlog growth, execution,
task generation, coordination, and low activity were derived from queue and
action variables. They are useful descriptive diagnostics, but are too close to
the manipulated variables to support attractor claims without residualization
and strong null controls.

\section{Experiment Timeline}

\begin{longtable}{p{0.16\linewidth}p{0.26\linewidth}p{0.46\linewidth}}
\hline
Stage & Question & High-level outcome \\
\hline
A0/A1 & Can OmegaSim produce deterministic reproducible run artifacts? &
Yes. Established configs, manifests, metrics, events, summaries, seeded runs,
and test discipline. No complex-dynamics claim. \\
\hline
A2 & Do attention policies and task pressure create trajectory regimes? &
Pressure robustly increases queue depth and queued age. This is a real system
effect, but mostly queue-flow/load arithmetic. \\
\hline
A2 scheduler ablation & Is the pressure effect caused by the quota scheduler? &
No. Random-available scheduling changes secondary curve shape but does not
remove the robust pressure response. \\
\hline
A2 exogenous arrivals & Can backlog pressure be induced independently of
agent task-creation pressure? & Yes. Exogenous arrivals reproduce backlog
pressure while agent-created tasks remain near control. \\
\hline
A3 & Can queue-flow and service-capacity accounting explain A2 effects? &
Yes. Demand pressure, finite service, work opportunity, completion fraction,
and queued age explain the robust effects. Lagged synchrony remains secondary. \\
\hline
A4 & Does two-hive transfer coupling induce meaningful cross-hive structure? &
Transfer coupling produces measurable completion synchronization, but the
signal is explained by queue-flow/service/action accounting controls. \\
\hline
A5 & Does resource-bounded prediction create residual structured dynamics? &
Forecast skill improves, but full-accounting residual structure and
compression gates fail. \\
\hline
A6/A6.1/A6.2 & Do logistic appraisal, artifact handoff, and longer-horizon
recurrence produce structured residual dynamics? & No promotion. Logistic
effects fail linear/source-preserving nulls and source-accounting checks. \\
\hline
A7 & Does a richer semantic/artifact field create residual semantic dynamics? &
No promotion. Field variation and source schema work, but null/residual gates
fail and productivity degrades. \\
\hline
A7.2 & Can delayed artifact-mediated endogenous prediction do better? &
Active next experiment. Preregistration and config/schema contract are frozen;
simulator/source-ledger mechanics and smoke results are next. \\
\hline
\end{longtable}

\section{Successful Results}

\subsection{Reproducible Experimental Harness}

OmegaSim now has a reproducible artifact discipline:

\begin{itemize}
\item deterministic YAML configurations;
\item seeded runs;
\item run manifests;
\item metric and event tables;
\item checked-in result summaries;
\item fail-closed automation guards;
\item paired-seed and null-control analyzers;
\item source-ledger reconstruction checks for later artifact-state experiments.
\end{itemize}

This is not merely software hygiene. It has shaped the science by repeatedly
preventing overinterpretation of attractive-looking but confounded signals.

\subsection{Robust Pressure and Queue-Flow Effects}

A2 found a real pressure response. In the extreme-pressure baseline over seeds
1--69, final queue depth for the internal-improvement policy changed as:

\begin{center}
\begin{tabular}{lrrr}
condition & normal & medium & extreme \\
\hline
final queue depth & 19.681159 & 45.956522 & 56.057971 \\
\end{tabular}
\end{center}

The high-minus-normal delta was 36.376812. This was the strongest early
effect, and it is robust as a queue/load phenomenon.

The scheduler ablation showed that the effect is not primarily an artifact of
the quota scheduler. In the seed 70--99 holdout:

\begin{center}
\begin{tabular}{lrrr}
scheduler & normal & medium & extreme \\
\hline
quota balance & 21.5 & 48.2 & 57.733333 \\
random available & 21.6 & 45.5 & 56.166667 \\
\end{tabular}
\end{center}

The random-available scheduler changes secondary metrics, but both schedulers
show the same qualitative pressure amplification.

A3 then converted this into a cleaner mechanism interpretation. Demand
pressure increases load-normalized backlog and queued age; service capacity
reduces backlog and improves completion fraction; exogenous arrivals reproduce
backlog pressure while agent-created task counts remain near control. The
scientific conclusion is that OmegaSim has a real queue-flow/service regime,
but this is not yet an attractor regime.

\subsection{Two-Hive Coupling Produces Measurable Synchronization Diagnostics}

A4 introduced a two-hive scaffold with no coupling, direct transfer, delayed
transfer, and shuffled transfer modes. In the seed 100--129 holdout, transfer
volume increased massively and deterministically for all transfer modes
relative to no coupling. Delayed coupling also produced completion-fraction
correlations that initially survived a circular-shift temporal null:

\begin{center}
\begin{tabular}{lrrr}
endpoint & observed minus null & outside null CI & positive rate \\
\hline
completion fraction corr lag 0 & 0.719486 & yes & 0.966667 \\
completion fraction corr lag 2 & 0.657053 & yes & 0.900000 \\
\end{tabular}
\end{center}

This was the most interesting A4 signal. However, the later accounting-control
analyzer showed that after combined non-tautological controls for clock,
queue/load, action mix, work opportunity, and transfer timing, the residual was
no longer outside the null interval:

\begin{center}
\begin{tabular}{lrrr}
endpoint & observed minus null & outside null CI & positive rate \\
\hline
completion fraction corr lag 0 & 0.165537 & no & 0.700000 \\
completion fraction corr lag 2 & 0.142740 & no & 0.666667 \\
\end{tabular}
\end{center}

This is a successful diagnostic result but a null mechanism result. Delayed
two-hive transfer does create synchronization-like traces, but the traces are
accounting-explainable in the current setup.

\subsection{Resource-Bounded Predictors Improve Forecast Skill}

A5 found a real forecast-skill manipulation. In the eight-condition
confirmatory run over seeds 7--16, predictive conditions improved forecast
skill under matched demand streams:

\begin{center}
\begin{tabular}{lr}
contrast & forecast-skill delta \\
\hline
linear minus reactive & 0.033517 \\
nonlinear minus reactive & 0.084033 \\
high-budget nonlinear minus reactive & 0.094796 \\
linear minus shuffled & 0.045425 \\
nonlinear minus nonlinear shuffled & 0.095941 \\
high-budget nonlinear minus high-budget shuffled & 0.106704 \\
\end{tabular}
\end{center}

This is the main positive A5 result: bounded predictors can improve forecast
skill in the simulator. But forecast skill did not translate into
promotion-worthy residual collective dynamics, as discussed below.

\subsection{Source-Ledger and Fail-Closed Analysis Infrastructure}

A6 and A7 were scientifically negative, but technically useful. They added or
validated source-accounting ideas needed for any serious future artifact-field
or semantic-field claim:

\begin{itemize}
\item artifact-update source fields;
\item reconstruction checks from source components;
\item source-preserving label shuffles;
\item timing-broken matched-count nulls;
\item residual recurrence analyzers;
\item semantic/artifact field completeness checks;
\item productivity guardrails.
\end{itemize}

A7.2 now builds directly on this infrastructure.

\section{Null and Fail-Closed Results}

\subsection{No Robust Evidence for Lobe Grammar from A2/A3}

The early lobe-like summaries are best interpreted as queue and action
diagnostics. A3 explicitly closed the A2 residual lobe-grammar hypothesis under
the current single-hive simulator and label scheme. Load, service capacity,
completion fraction, work opportunity, and queue age explain the observed
effects well enough that no independent lobe grammar is needed.

\subsection{Two-Hive A4 Does Not Yet Support Independent Cross-Hive Grammar}

The A4 delayed-coupling result was interesting, but too narrow. Completion
synchronization survived a simple temporal null but not the combined
non-tautological accounting controls. Also, the two-hive shuffled control is
degenerate for target-randomization purposes: with only two hives, each source
hive has only one non-source target. Therefore direct versus shuffled is mostly
a schema/opportunity check, not a real phase or target null.

This strongly motivates moving to three hives for serious phase and target
nulls.

\subsection{A5 Forecast Skill Did Not Become Residual Structure}

The A5 eight-condition closure is one of the clearest scientific boundaries.
Forecast skill improved, but the full-accounting residual-state predictability
contrasts stayed inside paired label-permutation intervals:

\begin{center}
\begin{tabular}{lrr}
contrast & residual mean delta & positive rate \\
\hline
linear minus reactive & 0.114949 & 0.5 \\
nonlinear minus reactive & 0.077593 & 0.7 \\
high-budget nonlinear minus reactive & 0.162115 & 0.6 \\
linear minus shuffled & -0.053977 & 0.5 \\
nonlinear minus nonlinear shuffled & -0.091333 & 0.3 \\
high-budget nonlinear minus high-budget shuffled & -0.006811 & 0.5 \\
\end{tabular}
\end{center}

The read-only residual-compression audit confirmed closure. Intermediate
predictors did not beat reactive, oracle smoothing, and timing-broken or
spend-only null expectations in a way that supports structured dynamics
language. The defensible A5 claim is narrow: resource-bounded prediction
improves forecast skill under matched single-hive demand streams.

\subsection{A6 Logistic Appraisal Did Not Pass Nulls}

A6 introduced logistic appraisal and artifact handoff ideas. The initial A6
smoke showed tiny endpoint differences, but logistic did not cleanly beat the
amplitude-matched linear condition. For example, in the seed 1--2 A6 analysis
gate:

\begin{center}
\begin{tabular}{lrrr}
contrast & artifact utility delta & queue-depth delta & completion delta \\
\hline
logistic vs linear & 0.001707 & 1.5 & -0.018187 \\
logistic vs phase-shuffled & 0.165515 & -2.0 & 0.059544 \\
logistic vs threshold-shuffled & 0.035120 & -1.5 & 0.061291 \\
\end{tabular}
\end{center}

The linear contrast had seed-sign disagreement, higher queue depth, and lower
completion fraction.

A6.1 added source-preserving nulls. All eight pilot null-gate rows were labeled
``null removes endpoint advantage''. Logistic readiness/utility advantages did
not survive source-label-shuffled or handoff-success timing-broken
matched-count controls.

\subsection{A6.2 Long-Horizon Validation Did Not Show Recurrence Advantage}

A6.2 extended the logistic-appraisal setup to 96 ticks. The analyzer computed
all recurrence rows, but closed conservatively:

\begin{verbatim}
overall analyzer status: conservative_closure
delta gate status:
  closure_no_recurrence_advantage: 98
  eligible_for_cross_seed_direction_check: 32
\end{verbatim}

Rows labeled eligible were not promotion evidence; they only passed local
same-seed checks against one control. The preregistered rule required the same
target variable to beat linear and source-preserving nulls with paired
cross-seed agreement. That did not happen. Dominant artifact-update sources
remained tied to handoff accounting or noise.

\subsection{A7 Semantic Field Did Not Pass Residual/Null Gates}

A7 added richer semantic/artifact fields and long-horizon validation. This
improved the semantic source-accounting path, but the logistic semantic-field
condition did not beat the required controls. In the A7 long-horizon validation
over seeds 1--2:

\begin{verbatim}
overall analyzer status: fail_closed_residual_null_gate
required field status: pass=12
source reconstruction status: pass=12
field variation status: pass=12
prediction/work-budget competition status: pass=12
residualization status: computed=72
null contrast status: computed=60
\end{verbatim}

Backlog-adjusted productivity for the logistic condition was worse than the
semantic-off baseline in both seeds:

\begin{center}
\begin{tabular}{lrr}
condition & seed 1 & seed 2 \\
\hline
logistic semantic coupling & 0.716667 & 0.654321 \\
semantic-off baseline & 1.543478 & 0.707865 \\
amplitude-matched linear & 0.689655 & 0.650000 \\
source-preserving label shuffle & 0.854167 & 0.628205 \\
phase shuffle & 0.661290 & 0.629630 \\
timing-broken null & 0.537313 & 0.511628 \\
\end{tabular}
\end{center}

No semantic/artifact target passed against all five controls in both seeds.
Several fields were still dominated by source-ledger components such as
semantic noise, artifact handoff, ambient decay, or clipping residuals.

\section{Scientific Interpretation So Far}

The evidence so far argues against a simple story in which pressure, task
queues, or shallow logistic thresholds automatically produce attractor-like
collective dynamics. The current conservative interpretation is:

\begin{enumerate}
\item OmegaSim has robust queue-flow/service dynamics.
\item It has measurable synchronization diagnostics under delayed coupling.
\item It can improve forecast skill with bounded predictors.
\item But none of these has yet produced residual structured dynamics that
survive accounting controls and nulls.
\end{enumerate}

This does not mean the larger hypothesis is unpromising. Rather, it sharpens
the likely mechanism requirements. Complex structured dynamics probably need
more than queue pressure and more than a predictor bolted onto a fixed
workflow. They likely require:

\begin{itemize}
\item nonlinear logistic or thresholded dependence of one agent's or hive's
activity on another's lagged state;
\item explicit delays, so prediction and artifact updates cannot help on the
same tick;
\item bounded prediction cost, so prediction competes with work;
\item imperfect anticipation, so prediction errors reshape the future state to
be predicted;
\item artifact fields that are not merely queue/action labels in disguise;
\item hysteresis, fatigue, and adaptive thresholds;
\item at least three hives for nondegenerate phase and target nulls;
\item promotion endpoints based on residual structure, not raw throughput.
\end{itemize}

This is exactly what A7.2 and the downstream three-hive ring are designed to
test.

\section{Active Next Experiment: A7.2}

\subsection{Scientific Question}

A7.2 asks whether delayed, artifact-mediated, resource-bounded endogenous
prediction creates residual prediction/artifact structure that survives full
accounting and strong null controls.

The intended mechanism is more endogenous than A5. Agents choose among:

\begin{verbatim}
predict
work
review
synthesize
\end{verbatim}

Action utilities are logistic functions of lagged forecast error, lagged
uncertainty, artifact readiness, contradiction, risk, fatigue, and adaptive
thresholds. Prediction consumes scarce work opportunity. Forecast and artifact
updates enter delayed queues and cannot affect action on the tick that created
them.

\subsection{Frozen Smoke Parameters}

The first A7.2 smoke is intentionally small:

\begin{verbatim}
seeds: 1, 2
horizon ticks: 48
forecast delay ticks: 2
artifact delay ticks: 3
prediction cost work fraction: 0.25
max prediction work fraction per tick: 0.40
fatigue decay: 0.20
fatigue increment predict: 0.08
fatigue increment work: 0.05
fatigue increment review: 0.04
fatigue increment synthesize: 0.06
threshold learning rate error: 0.05
threshold recovery rate: 0.02
threshold min/max: -2.0, 2.0
\end{verbatim}

These values may not be tuned after inspecting results. If they are wrong, the
experiment should close or be replaced by a new preregistration.

\subsection{A7.2 Conditions}

The first smoke must include:

\begin{verbatim}
zero_budget_reactive
intermediate_endogenous_delayed_prediction
high_budget_oracle_smoothing
amplitude_matched_linear_delayed_prediction
same_tick_logistic_prediction
phase_shuffled_lag_input
threshold_shuffled
source_preserving_artifact_label_shuffle
spend_only_replay
artifact_off_source_ledger_null
\end{verbatim}

Only the intermediate endogenous delayed-prediction condition is a candidate
positive condition. The oracle is a positive control, not a promotion target.

\subsection{A7.2 Primary Endpoints}

Promotion requires residual structure, not merely forecast skill:

\begin{verbatim}
forecast_skill_per_prediction_spend
full_accounting_residual_lag1_predictability
nearest_neighbor_residual_forecast_error
delay_embedding_recurrence_score
residual_transition_compressibility
lead_lag_mediation_error_to_spend_to_artifact_to_residual
source_ledger_reconstruction_status
\end{verbatim}

Residual endpoints must control for tick, demand phase, arrivals, service
capacity, action opportunity, work budget, prediction spend, remaining work
budget, backlog, queued age, completion fraction, starvation, spend
volatility, work-budget volatility, and source-ledger queue accounting.

\subsection{A7.2 Productivity Guardrails}

The candidate positive condition is ineligible for promotion if it causes
large degradation in completion fraction, backlog, queued age, starvation,
prediction-spend volatility, or work-budget volatility. The smoke guardrails
are deliberately modest and conservative.

\section{Downstream Three-Hive Ring}

Regardless of A7.2 outcome, the next phase is a separate three-hive ring
preregistration and bounded experiment family. This is not a rescue maneuver;
it is a scientific recognition that two hives are structurally inadequate for
meaningful target/phase nulls.

The proposed three hives should have biased roles or task tendencies, for
example:

\begin{itemize}
\item exploration/research;
\item formalization/implementation;
\item synthesis/review.
\end{itemize}

The ring should include delayed artifact transfer and cross-hive prediction
costs. Candidate conditions include:

\begin{verbatim}
no coupling
linear delayed coupling
same-tick logistic coupling
delayed logistic coupling
heterogeneous delay coupling
prediction-cost-feedback coupling
adaptive-threshold coupling
\end{verbatim}

Required controls should include target shuffles, phase shuffles,
transfer-opportunity matching, source-ledger reconstruction, spend-only replay,
and artifact-off or semantic-label controls where appropriate.

\section{External Feedback Questions}

The most useful external feedback would address the following questions:

\begin{enumerate}
\item Is A7.2 the right next single-hive mechanism, or should the project move
directly to three-hive dynamics?
\item Are the A7.2 nulls sufficient to prevent a false positive from backlog,
same-tick leakage, spend-only timing, oracle smoothing, or artifact labels?
\item What residual endpoint is most likely to detect ``structured but not
periodic'' dynamics in a short deterministic smoke?
\item Should the A7.2 artifact state remain low-dimensional, using readiness,
coherence, contradiction, and risk, or should a richer semantic activation
vector be introduced before the three-hive ring?
\item In the three-hive ring, what logistic dependency should be primary:
cross-hive prediction of demand, cross-hive artifact readiness, cross-hive
contradiction/risk, or cross-hive allocation pressure?
\item What would be a convincing null for ``structured grammatical patterns in
a strange attractor'' rather than ordinary phase-locked throughput?
\end{enumerate}

\section{Bottom Line}

OmegaSim has not yet elicited the desired structured strange-attractor
scenario. The positive results are real but bounded: robust pressure/service
effects, interpretable transfer/synchronization diagnostics, and genuine
forecast-skill improvements. The negative results are also informative:
apparent lobe, recurrence, semantic, and synchronization signals have so far
collapsed under accounting controls, source-preserving nulls, or productivity
guardrails.

The next best hypothesis is that complex dynamics require endogenous costly
prediction with delayed artifact mediation, and then a nondegenerate multi-hive
ring. A7.2 is the current test of the first half of that hypothesis; the
three-hive ring is the pre-authorized test of the second half.

\end{document}
